%%
%% spring-lecture03.tex
%% 
%% Made by Alex Nelson <pqnelson@gmail.com>
%% Login   <alex@lisp>
%% 
%% Started on  2026-04-04T10:52:42-0700
%% Last update 2026-04-04T10:52:42-0700
%% 

\lecture{}

\begin{definition}
Let $x=c_{0}+z\in A=\generalizedquat{a}{b}{\FF}$ be an element of the
generalized quaternions where $z=\quat{i}c_{1}+\quat{j}c_{2}+\quat{k}c_{3}$.
Let $y=d_{0}+w\in A$ where $w=\quat{i}d_{1}+\quat{j}d_{2}+\quat{k}d_{3}$.
We define the bilinear form $\beta$ on $A$ by
\begin{equation}
\begin{split}
\beta(x,y) &= \frac{1}{2}(\nu(x+y)-\nu(x)-\nu(y))\\
&=c_{0}d_{0}-ac_{1}d_{1}-bc_{2}d_{2}+abc_{3}d_{3}.
\end{split}
\end{equation}
We see that $\beta$ is symmetric [$\forall x,y\in A\ldotp\beta(x,y)=\beta(y,x)$] and
nondegenerate [$\forall x\in A\ldotp(\forall y\in A\ldotp\beta(x,y)=0)\implies x=0$].

Also observe on ``purely imaginary'' guys, we have $\beta(z,w)=\frac{-1}{2}(zw-wz)$ and $\nu(z)=-z^{2}$.
\end{definition}

\begin{note}
To see this, observe
\begin{subequations}
  \begin{align}
\beta(x,y) &= \frac{1}{2}((x+y)(x+y)^{*}-xx^{*}-yy^{*})\\
&=\frac{1}{2}(xy^{*}+yx^{*})\\
&=\frac{1}{2}((c_{0}+z)(d_{0}-w)+(d_{0}+w)(c_{0}-z))\\
&=c_{0}d_{0}-\frac{1}{2}(zw+wz)\\
&=c_{0}d_{0}-\frac{1}{2}(\quat{i}^{2}2c_{1}d_{1}+\quat{j}^{2}2c_{2}d_{2}+\quat{k}^{2}2c_{3}d_{3})\\
&=c_{0}d-{0}-ac_{1}d_{1}-bc_{2}d_{2}+abc_{3}d_{3}.
  \end{align}
\end{subequations}
Moreover, identifying $\FF$ with $1_{A}\FF=Z(A)$ allows us to relate
the norm with the bilinear form as $\nu(x)=\beta(x,x)$.
\end{note}

\begin{puzzle}
When are two generalized quaternion algebras isomorphic
$\generalizedquat{a}{b}{\FF}\iso\generalizedquat{a'}{b'}{\FF}$ as $\FF$-algebras?
\end{puzzle}

\begin{lemma}\label{lemma:1.12}
Let $A=\generalizedquat{a}{b}{\FF}$ and $A'=\generalizedquat{a'}{b'}{\FF}$
be equipped with norms $\nu$ and $\nu'$ respectively. Then $A\iso A'$
as $\FF$-algebras if and only if there exists a bijective $\FF$-linear
map $\phi\colon A_{+}\to A'_{+}$ such that for all $z\in A_{+}$ we
have $\nu'(\phi(z))=\nu(z)$.
\end{lemma}

\begin{proof}
\forwardproof\ For $x=c_{0}+z$,
\begin{subequations}
  \begin{align}
x^{2}&=c^{2}+z^{2}+z(2c)\\
&=(c^{2}-\nu(z))+z(2c)
  \end{align}
\end{subequations}
where $c^{2}-\nu(z)\in\FF$ and $z(2c)\in A_{+}$.
Then $x^{2}\in\FF$ if and only if $z=0$ or $c=0$; and
\begin{equation}\label{eq:spring-lec3:star:math120c}
(x\neq0\land x\in A_{+})\iff(x\notin Z(A)\land x^{2}\in Z(A)).
\end{equation}
Since every algebra isomorphism
\begin{equation}
\phi\colon A\to A'
\end{equation}
satisfies $\phi(Z(A))=Z(A')$, we have $\phi(A_{+})=A'_{+}$ by~\eqref{eq:spring-lec3:star:math120c}.
If $z\in A_{+}$, then since
\begin{equation}
\beta(z,w)=\frac{-1}{2}(zw+wz)\quad\mbox{and}\quad\nu(z)=-z^{2},
\end{equation}
we see
\begin{equation}
\nu'(\phi(z))=-\phi(x)^{2}=\phi(-z^{2})=\phi(\nu(z))=\nu(z)
\end{equation}
for all $z\in A_{+}$.
In other words, $\phi$ preserves the norm (hence proving the forward
direction).

\backwardproof\ Conversely, let $\phi\colon A\to A'$ be a bijective
linear map such that for all $z\in A_{+}$ we have $\nu'(\phi(z))=\nu(z)$.
To show that $A\iso A'$ as $\FF$-algebas, it sufices to check that
$\{1,\phi(\quat{i}),\phi(\quat{j}),\phi(\quat{k})\}$ satisfies the
defining relations of the generalized quaternion algebra.
The first three relations are easily checked
\begin{equation}
\phi(\quat{i})^{2}=-\nu'(\phi(\quat{i}))=\nu'(\quat{i})=a',
\end{equation}
and in a similar manner,
\begin{equation}
\phi(\quat{j})^{2}=b'.
\end{equation}
Then
\begin{subequations}
  \begin{align}
\phi(\quat{i})\phi(\quat{j})+\phi(\quat{j})\phi(\quat{i})
&=-2\beta'(\phi(\quat{i}),\phi(\quat{j}))\\
&=-2\beta(\quat{i},\quat{j})\\
&=\phi(\quat{i})\phi(\quat{j})+\phi(\quat{j})\phi(\quat{i})\\
&=0,
  \end{align}
\end{subequations}
and so
\begin{equation}
\phi(\quat{j})\bigl(\phi(\quat{i})\phi(\quat{j})\bigr)=-\phi(\quat{j})^{2}\phi(\quat{i})=-b'\phi(\quat{i}).
\end{equation}
On the other hand
\begin{equation}
\bigl(\phi(\quat{i})\phi(\quat{j})\bigr)\phi(\quat{j})=\phi(\quat{i})\phi(\quat{j})^{2}=b'\phi(\quat{i})\neq-b'\phi(\quat{i}),
\end{equation}
which implies $\phi(\quat{i})\phi(\quat{j})\notin Z(A')=\FF$, and
\begin{equation}
\bigl(\phi(\quat{i})\phi(\quat{j})\bigr)^{2}=-a'b'\in Z(A'),
\end{equation}
hence $\phi(\quat{i})\phi(\quat{j})\in A'_{+}$.

We want to show
$\{\phi(\quat{i}),\phi(\quat{j}),\phi(\quat{k})\}$ is a basis for
$A'_{+}$. Suffices to show they are linearly independent. If
\begin{equation}
0=z'=c_{1}\phi(\quat{i})+c_{2}\phi(\quat{j})+c_{3}\phi(\quat{i})\phi(\quat{j}),
\end{equation}
then
\begin{equation}
0=\phi(\quat{i})0=\phi(\quat{i})z'=\underbrace{a'c_{1}}_{\in\FF}+\underbrace{\phi(\quat{i})\phi(\quat{j})c_{2}+\phi(\quat{j})a'c_{3}}_{\in A'_{+}},
\end{equation}
so $c_{1}=0$ and similarly multiplying by $\phi(\quat{j})z'=0$ gives
$c_{2}=c_{3}=0$. Thus they are linearly independent and form a basis.

Then
\begin{equation}
\Psi\colon A\to A'
\end{equation}
such that $\Psi(1)=1$, $\Psi(\quat{i})=\phi(\quat{i})$,
$\Psi(\quat{j})=\phi(\quat{j})$, and
$\Psi(\quat{k})=\phi(\quat{i})\phi(\quat{j})$, then this is clearly an
isomorphism of $\FF$-algebras.
\end{proof}

\begin{corollary}\label{cor:1.13}
We have $A=\generalizedquat{a}{b}{\FF}\iso A'=\generalizedquat{a'}{b'}{\FF}$
if and only if the quadratic forms $ax_{1}^{2}+bx_{2}^{2}+abx_{3}^{2}$
and $a'x_{1}^{2}+b'_{1}x_{2}^{2}+a'b'x_{3}^{2}$ (i.e., there is an
invertible $3\times3$ matrix $\delta\in M_{3}(A_{+})$ such that $\diag(a,b,ab)=\delta^{\mathsf{T}}\diag(a',b',a'b')\delta$).
\end{corollary}

\begin{proof}
This follows from the matrix representations of $\nu=\diag(a,b,ab)$
and $\nu'=\diag(a',b',a'b')$ and $\phi=\delta$, since
\begin{equation}
\nu'(\phi(z))=(\delta\vec{c})^{\mathsf{T}}\diag(a',b',a'b')(\delta\vec{c})
\end{equation}
and
\begin{equation}
\nu(z)=\vec{c}^{\mathsf{T}}\diag(a,b,ab)\vec{c}
\end{equation}
and $\nu(z)=\nu'(\phi(z))$.
\end{proof}

\subsection{Modules}

\begin{convention}
Unless otherwise noted, $A$ is a nontrivial $R$-algebra. We fix $R$ as
well, and ``algebra'' means ``$R$-algebra''.

Here ``$A$-module'' means ``right $A$-module''.

``Ideal of $A$'' means two-sided ideal, and we denote it as $I\ideal A$.
\end{convention}

\begin{node}
Let $\theta\colon A\to B$ be an $R$-algebra morphism, let $M$ be a
$B$-module. We can define $A$-module structure on $M$ by
\begin{equation}
\forall u\in M\ldotp\forall x\in A\ldotp ux := u\theta(x).
\end{equation}
We denote this induced$A$-module structure on $M$ as $M_{A}$ or $MA$.

When $\theta\colon A\into B$ is a subalgebra, this induced $A$-module
structure morphism $M\to M_{A}$ is sometimes called the
\define{Forgetful Functor}.

When $\theta\colon A\to B$ and $B=A/I$ for some ideal $I\ideal A$, the
$A$-action on $M$ is $ux=u(x+I)$.
\end{node}

\begin{lemma}\label{lemma:2.1}
Let $\theta\colon A\to B$ be an $R$-algebra morphism. Let $M$, $N$ be
$B$-modules. Then the following all hold:
\begin{enumerate}
\item $(M\oplus N)_{A}\iso M_{A}\oplus N_{A}$;
\item $\hom_{B}(M,N)\subset\hom_{A}(M_{A},N_{A})$, and if $\theta$ is
  surjective then this is an equality;
\item If $N\subset M$ is a submodule, then $N_{A}\subset M_{A}$ is a
  submodule. If $\theta$ is surjective, and writing $S(M):=\{\mbox{submodules of }M\}$,
  then $S(M)=S(M_{A})$.
\end{enumerate}
\end{lemma}

\begin{proof}
Homework.
\end{proof}

\begin{definition}\label{defn:2.2}
Let $M$ be an $A$-module, let $X\subset M$ be a subset.
We define the \define{Annihilator} of $X$ to be the set
$\Annihilator(X)=\{a\in A\mid\forall x\in X\ldotp xa=0\}$.
\end{definition}

\begin{lemma}\label{lemma:2.3}
Let $M$, $M'$, $\{M_{i}\}_{i\in J}$ be $A$-modules.
Let $X\subset M$, $Y\subset M$ be subsets of $M$. Then the following
all hold:
\begin{enumerate}
\item $\Annihilator(X)$ is a right-ideal of $A$. If $X$ is a
  submodule, then $\Annihilator(X)\ideal A$.
\item If $X\subset Y$, then $\Annihilator(X)\supset\Annihilator(Y)$
\item If $M\iso M'$, then $\Annihilator(M)=\Annihilator(M')$
\item If $M=\sum_{i\in J}M_{i}$, then $\Annihilator(M)=\bigcap_{i\in J}\Annihilator(M_{i})$;
\item If $M$ is a right ideal of $A$, then $\Annihilator(A/M)$ is the
  largest ideal of $A$ contained in $M$.
\end{enumerate}
\end{lemma}

\begin{proof}
Items (1)--(4) are obvious.

(5) By (1), writing
\begin{equation}
X=A/M\subset A/M,
\end{equation}
we find $\Annihilator(A/M)\properideal A$ and since
\begin{equation}
\Annihilator(A/M)=\{x\in A\mid Ax\subset M\}
\end{equation}
and $1_{A}\in A$, so $\Annihilator(A/M)\subset M$. If $K\ideal A$ is
another ideal such that $K\subset M$, then $AK\subset M$ which means $K\subset\Annihilator(A/M)$.
\end{proof}

\begin{proposition}\label{prop:2.4}
For any surjective algebra morphism $\theta\colon A\to B$ and for any
$A$-module $N$, then the following are equivalent:
\begin{enumerate}
\item there exists a $B$-module $M$ such that $N=M_{A}$,
\item $\ker(\theta)\subset\Annihilator(N)$.
\end{enumerate}
\end{proposition}

\begin{proof}
\forwardproof\ Obvious.

\backwardproof\ Since $\theta$ is surjective, for any $y\in B$ there
exists some $x\in A$ such that $y=\theta(x)$, then $uy:=u\theta(x)=ux$
defines a $B$-module on $N=:M$, so $M_{A}=N$.
\end{proof}

\begin{definition}
Let $M$ be an $A$-module. We say $M$ is \define{Faithful} if $\Annihilator(M)=0$.
\end{definition}

\begin{corollary}
For an ideal $I$ of $A$ and an $(A/I)$-module $N$, we have $N$ is
faithful (i.e., $\Annihilator(N)=0$) if and only if $\Annihilator(N_{A})=I$.
\end{corollary}